%% LobsterCXR Paper — MICCAI 2027 Submission
%% Compile: pdflatex lobstercxr_paper.tex (×2)
%% Pre-rendered figures: figures/architecture.png, figures/ddx_chain.png
%% Generated from Mermaid source: figures/*.mmd

\documentclass[runningheads,10pt]{llncs}

\usepackage{amsmath,amssymb,amsfonts}
\usepackage{graphicx}
\usepackage{booktabs}
\usepackage{multirow}
\usepackage{array}
\usepackage{xcolor}
\usepackage{colortbl}
\usepackage{url}
\usepackage{algorithm}
\usepackage{algpseudocode}
\usepackage[colorlinks=true,linkcolor=blue,citecolor=blue,urlcolor=blue]{hyperref}

%% ------------------------------------------------------------
%% Custom notation
%% ------------------------------------------------------------
\newcommand{\Fv}{\mathbf{F}_v}
\newcommand{\yhat}{\hat{\mathbf{y}}}
\newcommand{\Sstruct}{S_{\text{struct}}}
\newcommand{\Sread}{S_{\text{read}}}
\newcommand{\Steach}{S_{\text{teach}}}
\newcommand{\Stotal}{S_{\text{disc}}}
\newcommand{\Rtotal}{R_{\text{total}}}
\newcommand{\Rteach}{R_{\text{teaching}}}
\newcommand{\Rclin}{R_{\text{clinical}}}
\newcommand{\Lce}{\mathcal{L}_{\text{ce}}}
\newcommand{\Lvlp}{\mathcal{L}_{\text{vlp}}}
\newcommand{\Lpair}{\mathcal{L}_{\text{pair\_margin}}}
\newcommand{\Lteach}{\mathcal{L}_{\text{teach}}}
\newcommand{\Ltotal}{\mathcal{L}_{\text{total}}}
\newcommand{\Ksep}{\mathcal{K}_{\text{sep}}}
\newcommand{\Cout}{\mathcal{C}_{\text{out}}}

%% Box for key formula
\newsavebox{\keybox}
\newcommand{\keyformula}[2][eq:key]{%
  \vspace{2mm}\noindent
  \fbox{\parbox{0.92\textwidth}{\centering\vspace{2mm}%
    \begin{math}\displaystyle\begin{aligned} #2 \end{aligned}\end{math}%
  \vspace{2mm}}}\vspace{2mm}\label{#1}}

%% Column types
\newcolumntype{C}{>{\centering\arraybackslash}m{2.1cm}}
\newcolumntype{L}{>{\raggedright\arraybackslash}m{3.2cm}}

\begin{document}

\title{LobsterCXR: A Multi-Agent Collaborative Framework for Differential Diagnosis and Educational Report Generation}

\author{First Author\textsuperscript{1} \and Second Author\textsuperscript{1,2}}
\institute{
University Name, City, Country \textsuperscript{1}
\and
Hospital Name, City, Country \textsuperscript{2}
\email{\{author1, author2\}@institution.edu}
}

\maketitle

\begin{abstract}
Automated chest X-ray (CXR) report generation has advanced considerably, yet existing methods lack three critical capabilities: (1) an explicit differential diagnosis reasoning loop for distinguishing confusable diseases; (2) evaluation systems that incorporate pedagogical readability alongside clinical accuracy; and (3) a unified multi-agent framework simultaneously addressing both reasoning and education.

Our framework, \textbf{LobsterCXR}, teams up four cognitively-role-driven agents: \textbf{Vision Analyst} (CheXagent ViT-L/14 encoding with multi-scale features and 14-class CheXpert confidence calibration), \textbf{Draft Writer} (Vicuna-7B + LoRA generation with Q-Former cross-modal compression and retrieval augmentation), \textbf{DDx Critic} (a four-stage ``Discuss-Differentiate'' reasoning chain producing structured comparative rationales for confusable diagnostic pairs), and \textbf{Teaching Evaluator} (three-dimensional pedagogical scoring across structure, readability, and teaching suitability). A dual-layer consensus mechanism coordinates inter-agent outputs via confidence gating and weighted fusion, with end-to-end optimization combining language modeling, contrastive discrimination, and GRPO reinforcement learning.

Experiments on MIMIC-CXR and IU X-Ray demonstrate that LobsterCXR matches or exceeds state-of-the-art clinical accuracy (CheXbert F1, RadGraph F1, GREEN) while significantly outperforming baselines on pedagogical readability metrics. Ablation studies confirm the marginal contribution of each agent and the consensus module.

\keywords{Chest X-ray; Report Generation; Multi-Agent Systems; Differential Diagnosis; Medical Education}
\end{abstract}

%% ============================================================
%% 1. INTRODUCTION
%% ============================================================
\section{Introduction}

Chest X-ray (CXR) is the most widely performed radiological examination globally, with over 1.4 billion procedures conducted annually~\cite{chexagent}. The global shortage of radiologists leads to prolonged report turnaround times and increased diagnostic error rates~\cite{tanno2023}. Automated CXR report generation has progressed from early RNN-based models~\cite{r2gen} to large vision-language models~\cite{radialog,chexagent,cxrmate2}, yet through systematic analysis of over 20 SOTA works we identify three critical gaps:

\textbf{Gap 1: Absence of differential diagnosis reasoning loops.} Every existing method~\cite{r2gen,chexagent,riha,cxrmate2} maps images straight to text in one pass---no ``generate-criticize-revise'' loop in sight. MARL-Rad~\cite{marlrad} adopts region-based multi-agent design but divides by anatomical regions rather than cognitive roles. No prior work integrates differential diagnosis critique as an optimizable reasoning loop within report generation.

\textbf{Gap 2: Pedagogical assessment absent from end-to-end optimization.} All standard metrics (RadGraph F1~\cite{radgraph}, CheXbert F1, GREEN~\cite{green}) focus exclusively on clinical completeness. MAARTA~\cite{maarta} and IMACT-CXR~\cite{imactcxr} serve educational purposes but require eye-tracking hardware or interactive input, precluding automated deployment.

\textbf{Gap 3: Lack of a unified framework.} No existing work simultaneously covers multi-role agent collaboration, differential diagnosis reasoning, and pedagogical readability optimization.

LobsterCXR tackles these gaps head-on. Concretely:
\begin{enumerate}
    \item \textbf{A cognitive-role-driven multi-agent architecture} for CXR report generation (\S3.1)---instead of splitting by anatomy, we split by thinking stage.
    \item \textbf{A discuss-consensus reasoning loop} (\S3.4) implementing explicit comparative differential diagnosis via a four-stage ``Extract$\to$Build$\to$Discriminate$\to$Output'' chain.
    \item \textbf{Pedagogical readability as an end-to-end optimizable objective} (\S3.5), introducing a three-dimensional teaching quality scoring system that participates in optimization through GRPO rewards, differentiable auxiliary losses, and inference gating.
    \item \textbf{A dual-layer consensus mechanism} (\S3.6) combining confidence gating and weighted fusion to manage inter-agent disagreement and enable iterative report refinement.
\end{enumerate}

%% ============================================================
%% 2. RELATED WORK
%% ============================================================
\section{Related Work}

\subsection{CXR Report Generation}
R2Gen~\cite{r2gen} combines DenseNet-121 with memory-driven Transformer decoders. RaDialog~\cite{radialog} introduces Vicuna-7B as decoder with LoRA fine-tuning. CheXagent~\cite{chexagent} trains a multimodal LLM on 28 CXR datasets. CXRMate-2~\cite{cxrmate2} introduces GRPO reinforcement learning. RIHA~\cite{riha} employs three-level optimal transport cross-modal alignment. Wu et al.~\cite{diseaseaware} propose disease-aware semantic tokens (DASTs) and dual-modal retrieval. All adopt single-model encoder-decoder architectures without inter-agent discussion or explicit differential diagnosis reasoning.

\subsection{Multi-Agent Medical AI}
AutoGen~\cite{autogen} provides general multi-agent conversation orchestration. MARL-Rad~\cite{marlrad} applies multi-agent RL to CXR report generation with anatomy-specific agents (left lung, right lung, heart) and a global integrator. XrayClaw~\cite{xrayclaw} proposes cooperative-competitive agents for CXR classification using competitive preference optimization. These works validate agent-based architectures for medical tasks but target classification rather than report generation with differential diagnosis.

\subsection{Educational Assessment Systems}
MAARTA~\cite{maarta} (MICCAI~2025) is the first multi-agent radiology teaching assistant, identifying perceptual errors via gaze pattern comparison between experts and trainees. IMACT-CXR~\cite{imactcxr} (ISBI~2026) proposes an interactive multi-agent CXR tutoring system with Bayesian Knowledge Tracing. Both advance the educational dimension but require hardware or real-time interaction, making them unsuitable for automated educational report generation.

\subsection{Our Positioning}
No prior work simultaneously addresses differential diagnosis reasoning loops, pedagogical readability as an optimization target, and multi-agent collaboration for CXR report generation. LobsterCXR aims to fill all three gaps within a unified framework.

%% ============================================================
%% 3. METHOD
%% ============================================================
\section{Method}

\subsection{Overall Architecture}

LobsterCXR mirrors how radiologists actually work: four cognitive stages---\textbf{Observe} (scanning image, localizing abnormalities), \textbf{Write} (generating initial draft), \textbf{Verify} (validating diagnoses, considering differentials), and \textbf{Teach} (structuring output for learner comprehension)---each assigned to a dedicated agent.

\begin{figure}[t]
\centering
\includegraphics[height=0.6\textheight]{figures/architecture.png}
\caption{LobsterCXR overall architecture. Four cognitively-role-driven agents collaborate through a dual-layer consensus module: Vision Analyst encodes CXR into multi-scale features and disease confidences; Draft Writer generates the initial report; DDx Critic conducts differential diagnosis analysis; Teaching Evaluator scores pedagogical quality. The consensus module determines whether to trigger revision based on divergence ($\tau_{\text{gate}}$) and teaching scores ($\mathbf{S}$).}
\label{fig:arch}
\end{figure}

\subsection{Agent A: Vision Analyst}

\textbf{What it does.} Encodes the input CXR at multiple granularities so downstream agents have both broad context and fine detail.

\textbf{Architecture.} The visual encoder is CheXagent's ViT-L/14~\cite{chexagent}, pre-trained on 28 CXR datasets. To capture findings at different scales, we add RIHA's Visual Feature Pyramid (VFP)~\cite{riha} with three resolution levels ($16\times/32\times/64\times$). A lightweight classification head maps the [CLS] token to CheXpert's 14 disease classes with a learned temperature-scaling calibration.

Let input image be $\mathbf{x} \in \mathbb{R}^{3 \times H \times W}$ with $H=W=224$:
\begin{equation}
\Fv, \mathbf{c} = \text{ViT}(\mathbf{x}) \label{eq:vit}
\end{equation}
where $\Fv \in \mathbb{R}^{L \times d}$ ($L=196$, $d=768$) are patch-level features and $\mathbf{c} \in \mathbb{R}^{d}$ is the [CLS] token. Multi-scale features:
\begin{equation}
\Fv^{\text{pyr}} = \{ \Fv^{(16)}, \Fv^{(32)}, \Fv^{(64)} \} = \text{VFP}(\Fv) \label{eq:pyramid}
\end{equation}
Calibrated 14-class predictions:
\begin{equation}
\yhat_A = \sigma\!\left( \frac{\mathbf{W}_c \mathbf{c} + \mathbf{b}_c}{\tau} \right) \in [0, 1]^{14} \label{eq:classifier}
\end{equation}

\subsection{Agent B: Draft Writer}

\textbf{Role.} Generates the initial structured report (Findings $\to$ Impression) with disease-aware attention and retrieval augmentation.

\textbf{Backbone.} Vicuna-7B + LoRA ($r=16$)~\cite{radialog}. Q-Former (32 learnable queries)~\cite{chexagent} for cross-modal compression. Disease-Aware Semantic Tokens (DASTs) and Dual-Modal Similarity Retrieval (DMSR)~\cite{diseaseaware}.

Q-Former compressed features:
\begin{equation}
\Fv^{\text{cond}} = \text{CrossAttn}(\mathbf{Q}, \Fv, \Fv) \in \mathbb{R}^{32 \times d} \label{eq:qformer}
\end{equation}
Disease-aware tokens from $\yhat_A$:
\begin{equation}
\mathbf{D}_{\text{tok}} = \mathbf{W}_d \yhat_A + \mathbf{b}_d \in \mathbb{R}^{14 \times d} \label{eq:dast}
\end{equation}
Retrieval tokens from database $\mathcal{H}$:
\begin{equation}
\mathbf{R}_{\text{tok}} = \text{DMSR}(\Fv, \mathcal{H}) \in \mathbb{R}^{8 \times d} \label{eq:dmsr}
\end{equation}
Autoregressive decoding with $\mathbf{Z}_0 = [\Fv^{\text{cond}}; \mathbf{D}_{\text{tok}}; \mathbf{R}_{\text{tok}}]$:
\begin{equation}
P(R_0 \mid \mathbf{x}) = \prod_{t=1}^{T} P_{\theta_B}(w_t \mid w_{<t}, \mathbf{Z}_0) \label{eq:gen}
\end{equation}

\subsection{Agent C: DDx Critic (Core Innovation)}

\textbf{Role.} Critically reviews the draft $R_0$ for differential diagnosis quality via a ``Discuss-Differentiate'' reasoning chain.

\textbf{Backbone.} Vicuna-13B + LoRA ($r=16$), independent adapter for role specialization. Label space interface: CheXbert~\cite{chexpert} for standardized entity mapping.

\textbf{Stage 1---Extract Findings.} Parse $R_0$ via CheXbert:
\begin{equation}
\mathcal{F} = \{ (a_i, o_i, s_i) \}_{i=1}^{N} = \text{CheXbert}(R_0) \label{eq:extract}
\end{equation}

\textbf{Stage 2---Build Differential Space.} Query disease-finding knowledge graph $\mathcal{K}$:
\begin{equation}
\mathcal{D}_i = \mathcal{K}.\text{query}(f_i) = \{ (d_j, \text{evidence}_j, w_j) \} \label{eq:diffspace}
\end{equation}

\textbf{Stage 3---Discriminate (Core).} For each confusable pair $(d_p, d_q)$, cross-modal contrastive reasoning:
\begin{equation}
\boxed{
\Stotal(d_p, d_q) = \frac{
    \sum_{e \in \Ksep(d_p, d_q)} \mathbf{1}[e \in \mathcal{F}] \cdot r(e, d_p)
    + \alpha \cdot v(d_p)
}{
    \sum_{e \in \Ksep(d_p, d_q)} \mathbf{1}[e \in \mathcal{F}] \cdot (r(e, d_p) + r(e, d_q))
    + \alpha \cdot (v(d_p) + v(d_q))
}
} \label{eq:disc}
\end{equation}
where $\Ksep$ are discriminating features between $d_p$ and $d_q$, $r(e,d)$ indicates feature-disease support, $v(d)=\yhat_A[d]$ is visual confidence, $\alpha$ weights visual evidence.

\textbf{Stage 4---Structured Output:}
\begin{equation}
\Cout = \{\text{claimed\_dx},\ \text{pairs}: (d_p,d_q,\Stotal,\text{rationale}),\ \text{missed},\ \text{overall}\} \label{eq:ddxout}
\end{equation}

\textbf{Key insight:} Eq.~\eqref{eq:disc} performs comparative reasoning over diagnosis \emph{pairs}, not independent per-disease logits. This enables explicit multi-evidence chain reasoning and competitive diagnosis differentiation---fundamentally different from all existing multi-label classification approaches.

\begin{figure}[t]
\centering
\includegraphics[width=\textwidth]{figures/ddx_chain.png}
\caption{DDx Critic's four-stage ``Discuss-Differentiate'' reasoning chain. Stage 3 implements contrastive scoring (Eq.~\eqref{eq:disc}) for explicit evidence-based discrimination between confusable diagnostic pairs.}
\label{fig:ddx}
\end{figure}

\subsection{Agent D: Teaching Evaluator}

\textbf{Role.} Non-generative scoring model quantifying pedagogical quality across three dimensions.

\textbf{Backbone.} PubmedBERT-base + 3 independent linear regression heads. RadGraph entity coverage~\cite{green}.

Three-dimensional scores $\mathbf{S} \in [0,1]^3 = (\Sstruct, \Sread, \Steach)$:
\begin{align}
\Sstruct &= 0.25 \cdot \mathbb{1}[\text{has\_headers}] + 0.30 \cdot \text{region\_coverage} \nonumber \\
&\quad + 0.15 \cdot \text{format\_consistency} + 0.30 \cdot \text{entity\_coverage} \label{eq:sstruct}\\[2mm]
\Sread &= 0.20 \cdot \ell_{\text{len}} + 0.30 \cdot \ell_{\text{density}} + 0.25 \cdot \ell_{\text{explain}} + 0.25 \cdot \ell_{\text{complexity}} \label{eq:sread}\\[2mm]
\Steach &= 0.30 \cdot \mathbb{1}[\text{has\_ddx}] + 0.25 \cdot \text{causal\_score} + 0.15 \cdot \text{neg\_findings} \nonumber \\
&\quad + 0.15 \cdot \mathbb{1}[\text{has\_rec}] + 0.15 \cdot \text{context\_score} \label{eq:steach}
\end{align}

\textbf{Optimization participation:} (a) GRPO reward: $\Rteach = \mathbf{w}^\top \mathbf{S}$ with $\mathbf{w}=[0.25,0.35,0.40]$; (b) differentiable auxiliary loss: $\Lteach = \text{MSE}(\mathbf{S}_{\text{pred}}, \mathbf{S}_{\text{target}})$; (c) inference gating: stop iteration when $\min(\mathbf{S}) \geq \eta_{\text{teach}}$.

\subsection{Consensus Mechanism}

\textbf{Layer 1---Confidence Gate.} Divergence between Draft Writer and DDx Critic for disease $d$:
\begin{equation}
\Delta(d) = |P_B(d) - P_C(d)| \label{eq:divergence}
\end{equation}
\begin{equation}
\text{Gate} = \begin{cases}
\text{trigger\_revision}, & \max_d \Delta(d) > \tau_{\text{gate}} \\
\text{skip\_to\_fusion}, & \text{otherwise}
\end{cases} \label{eq:gate}
\end{equation}

\textbf{Layer 2---Weighted Fusion.}
\begin{equation}
P_{\text{final}}(d) = \frac{w_B \cdot P_B(d) + w_C \cdot P_C(d) + w_{\text{con}} \cdot P_{\text{con}}(d)}{w_B + w_C + w_{\text{con}}} \label{eq:fusion}
\end{equation}
with consensus indicator:
\begin{equation}
P_{\text{con}}(d) = \begin{cases}
P_B(d), & \Steach \geq \eta_{\text{teach}} \\
P_C(d), & \Steach < \eta_{\text{teach}} \land \Delta(d) > \tau_{\text{gate}} \\
\frac{P_B(d) + P_C(d)}{2}, & \text{otherwise}
\end{cases} \label{eq:conind}
\end{equation}

\textbf{Stop condition:}
\begin{equation}
\text{Stop} \iff \bigl[ \max_d \Delta(d) < \tau_{\text{gate}} \land \min(\mathbf{S}) \geq \eta_{\text{teach}} \bigr] \lor [\text{rounds} \geq 3] \lor [\text{ROUGE-L} \geq 0.95] \label{eq:stop}
\end{equation}

\subsection{Training Paradigm and Loss Functions}

Three-stage progressive training.

\textbf{Stage 1: Agent A+B pre-training.}
\begin{equation}
\mathcal{L}_{\text{stage1}} = \Lce + \lambda_{\text{vlp}} \Lvlp \label{eq:stage1}
\end{equation}

\textbf{Stage 2: Agent C specialization.}
\begin{equation}
\mathcal{L}_{\text{stage2}} = \Lce^{(C)} + \lambda_{\text{pair}} \max(0, \gamma - (P_{\text{gt}}(d_1) - P_{\text{gt}}(d_2))) \label{eq:stage2}
\end{equation}

\textbf{Stage 3: Joint fine-tuning (GRPO).}
\begin{equation}
\mathcal{L}_{\text{stage3}} = -\mathbb{E}_{R \sim \pi_{\theta}} [\Rtotal(R)] + \beta_{\text{KL}} \cdot \text{KL}(\pi_{\theta} \| \pi_{\text{ref}}) \label{eq:stage3}
\end{equation}

Reward function:
\begin{equation}
\Rtotal(R) = \alpha \cdot \Rclin(R; \mathbf{y}_{\text{ref}}) + (1-\alpha) \cdot \Rteach(\mathbf{S}) \label{eq:reward}
\end{equation}

\textbf{Total loss:}
\keyformula[eq:total]{
\Ltotal = &\beta_1 \Lce + \beta_2 \Lvlp + \beta_3 \Lpair \\
&+ \beta_4 \Lteach - \beta_5 \mathbb{E}[\Rtotal] + \beta_6 \text{KL}(\pi_{\theta} \| \pi_{\text{ref}})
}

Coefficient schedule:
\begin{table}[t]
\centering
\caption{Stage-dependent loss coefficients.}
\label{tab:coeff}
\begin{tabular}{lcccccc}
\toprule
Stage & $\beta_1$ & $\beta_2$ & $\beta_3$ & $\beta_4$ & $\beta_5$ & $\beta_6$ \\
\midrule
Stage 1 & 1.0 & 0.5 & 0 & 0 & 0 & 0 \\
Stage 2 & 0 & 0 & 1.0 & 0 & 0 & 0 \\
Stage 3 & 0.5 & 0.1 & 0.5 & 1.0 & 0.01 & 0.05 \\
\bottomrule
\end{tabular}
\end{table}

%% ============================================================
%% 4. EXPERIMENTS
%% ============================================================
\section{Experiments}

\subsection{Datasets}

\begin{table}[t]
\centering
\caption{Datasets used in evaluation.}
\label{tab:datasets}
\begin{tabular}{llrl}
\toprule
Dataset & Purpose & Size & Source \\
\midrule
MIMIC-CXR~\cite{mimiccxr} & Train + Test & $\sim$190K~/~$\sim$3K & Beth Israel Deaconess \\
IU X-Ray & Cross-domain test & $\sim$5.2K~/~$\sim$2K & Indiana University \\
CheXpert Plus~\cite{chexpert} & DDx pre-training & $\sim$200K & Stanford Hospital \\
\bottomrule
\end{tabular}
\end{table}

\subsection{Baselines and Metrics}

\textbf{Baselines:} R2Gen~\cite{r2gen}, RaDialog~\cite{radialog}, MARL-Rad~\cite{marlrad}, CXRMate-2~\cite{cxrmate2}. Ablation variants: Ours$-$DDx, Ours$-$Teaching, Ours$-$Consensus, Ours (1-round).

\textbf{Clinical metrics:} CheXbert F1 (macro/micro), RadGraph F1~\cite{radgraph}, GREEN~\cite{green}, RadCliQ.
\textbf{Pedagogical metrics:} $\Sstruct$, $\Sread$, $\Steach$ (ours); Human Teaching Rating (1--5).
\textbf{Utility metrics:} Radiologist preference (2AFC), resident time saving.

\subsection{Implementation Details}

Hardware: 1$\times$ NVIDIA A100 (80GB). Code: \texttt{./experiments/}. Optimizer: AdamW.

\noindent\textbf{Agent A:} CheXagent ViT-L/14. VFP 3-level. 2-layer MLP (768$\to$256$\to$14). $\tau$ init 1.0.
\textbf{Agent B:} Vicuna-7B v1.5, LoRA $r=16$, $\alpha=32$. Q-Former 32 queries. DMSR Top-8.
\textbf{Agent C:} Vicuna-13B v1.5, LoRA $r=16$. CheXbert frozen. $\mathcal{K}$ from CheXpert+RadGraph.
\textbf{Agent D:} PubmedBERT-base (110M), 3-head regression, dropout 0.1, lr 2e-5.

\noindent\textbf{Training:} Stage 1: lr 5e-5, bs 16. Stage 2: lr 3e-5, bs 16. Stage 3: lr 1e-5, bs 8, group 4, $\beta_{\text{KL}}=0.05$, $\alpha=0.7$. Consensus: $\tau_0=0.15$, $\eta=0.1$, $\eta_{\text{teach}}=0.7$.

\subsection{Main Results}

\begin{table}[t]
\centering
\caption{Clinical accuracy comparison on MIMIC-CXR test set.}
\label{tab:main}
\small
\begin{tabular}{lCCCCC}
\toprule
Method & CheXbert (macro) & CheXbert (micro) & RadGraph F1 & GREEN & RadCliQ \\
\midrule
R2Gen~\cite{r2gen} & -- & -- & -- & -- & -- \\
RaDialog~\cite{radialog} & -- & -- & -- & -- & -- \\
MARL-Rad~\cite{marlrad} & -- & -- & -- & -- & -- \\
CXRMate-2~\cite{cxrmate2} & -- & -- & -- & -- & -- \\
\midrule
\textbf{LobsterCXR} & -- & -- & -- & -- & -- \\
\bottomrule
\end{tabular}
\end{table}

\begin{table}[t]
\centering
\caption{Teaching readability comparison.}
\label{tab:teach}
\small
\begin{tabular}{lCCCC}
\toprule
Method & $\Sstruct$ & $\Sread$ & $\Steach$ & Human Rating (1--5) \\
\midrule
R2Gen & -- & -- & -- & -- \\
RaDialog & -- & -- & -- & -- \\
MARL-Rad & -- & -- & -- & -- \\
CXRMate-2 & -- & -- & -- & -- \\
\midrule
\textbf{LobsterCXR} & -- & -- & -- & -- \\
\bottomrule
\end{tabular}
\end{table}

\subsection{Ablation Study}

\begin{table}[t]
\centering
\caption{Ablation: marginal contributions of each module (MIMIC-CXR).}
\label{tab:ablation}
\small
\begin{tabular}{lCCCC}
\toprule
Variant & CheXbert (macro) & GREEN & $\Steach$ & Avg. Rounds \\
\midrule
Full model & -- & -- & -- & -- \\
$-$\texttt{DDx} & -- & -- & -- & -- \\
$-$\texttt{Teaching} & -- & -- & -- & -- \\
$-$\texttt{Consensus} & -- & -- & -- & -- \\
1-round only & -- & -- & -- & 1.0 \\
\bottomrule
\end{tabular}
\end{table}

\subsection{Cross-Domain Generalization}

\begin{table}[t]
\centering
\caption{IU X-Ray cross-domain results (zero-shot transfer from MIMIC-CXR).}
\label{tab:cross}
\small
\begin{tabular}{lCCC}
\toprule
Method & CheXbert (macro) & $\Delta$ from MIMIC & $\Steach$ \\
\midrule
MARL-Rad & -- & -- & -- \\
CXRMate-2 & -- & -- & -- \\
\textbf{LobsterCXR} & -- & -- & -- \\
\bottomrule
\end{tabular}
\end{table}

\subsection{Blinded Evaluation}

\begin{table}[t]
\centering
\caption{Blinded evaluation: radiologist and medical student preferences (2AFC).}
\label{tab:blind}
\small
\begin{tabular}{lCCC}
\toprule
Comparison Pair & Radiologist Pref. & Student Pref. & Cohen's $\kappa$ \\
\midrule
LobsterCXR vs. CXRMate-2 & -- & -- & -- \\
LobsterCXR vs. MARL-Rad & -- & -- & -- \\
\bottomrule
\end{tabular}
\end{table}

\subsection{Computational Cost}

\begin{table}[t]
\centering
\caption{Computational cost comparison.}
\label{tab:cost}
\small
\begin{tabular}{lCCC}
\toprule
Method & Training (GPU h) & Inference / sample & Parameters \\
\midrule
R2Gen & -- & -- & 120M \\
RaDialog & -- & -- & $\sim$7B + LoRA \\
MARL-Rad & -- & -- & $\sim$7B + LoRA \\
CXRMate-2 & -- & -- & $\sim$7B + LoRA \\
\textbf{LobsterCXR} & -- & -- & $\sim$7B+13B + LoRA \\
\bottomrule
\end{tabular}
\end{table}

%% ============================================================
%% 5. DISCUSSION
%% ============================================================
\section{Discussion}

\subsection{Response to Three Gaps}

How well does LobsterCXR plug the three holes? Let's walk through each. The DDx Critic's four-stage chain (Eq.~\eqref{eq:extract}--\eqref{eq:ddxout}) turns differential diagnosis from a black-box end-to-end mapping into explicit, evidence-based pairwise comparisons (Gap~1). The Teaching Evaluator's three-dimensional scoring (Eq.~\eqref{eq:sstruct}--\eqref{eq:steach}) is, to our knowledge, the first attempt to turn pedagogical readability into a differentiable, end-to-end optimizable signal (Gap~2). The cognitive-role-driven agent division provides a unified framework addressing both diagnosis and education (Gap~3).

\subsection{Comparison with Prior Work}

\textbf{vs. MARL-Rad~\cite{marlrad}:} Their agents divide by anatomy (left lung, right lung, heart); ours divide by cognitive role. Each of our agents processes the full visual field but handles a distinct reasoning stage, enabling deeper per-stage specialization.

\textbf{vs. XrayClaw~\cite{xrayclaw}:} Their CPO validates multi-agent verification in classification. We extend this paradigm to free-text report generation with structured output constraints and a DDx-specific reasoning chain.

\textbf{vs. MAARTA~\cite{maarta}:} Their gaze-based perceptual error analysis is more granular but requires eye-tracking hardware. Our text-only approach has zero deployment overhead, enabling scalable educational report generation.

\textbf{vs. CXRMate-2~\cite{cxrmate2}:} Their GRPO operates on a single agent. Our framework extends RL-based optimization with multi-agent collaboration and multi-dimensional (clinical + teaching) rewards.

\subsection{Limitations}

(1) Our DDx knowledge graph is only as good as CheXpert and RadGraph's labels. Rare or emerging diseases will slip through. (2) The Teaching Evaluator scores via regression, which is less semantically flexible than prompting an LLM---but it backpropagates gradients, which we need for end-to-end training. (3) Three training stages are cumbersome. A single-stage end-to-end procedure would be cleaner, and we plan to move in that direction.

%% ============================================================
%% 6. CONCLUSION
%% ============================================================
\section{Conclusion}

We presented LobsterCXR, a multi-agent collaborative framework for CXR differential diagnosis and educational report generation. Four cognitive-role-driven agents (Vision Analyst, Draft Writer, DDx Critic, Teaching Evaluator) coordinate through a dual-layer consensus mechanism combining confidence gating and weighted fusion. The core innovations include: (1) a discuss-consensus reasoning loop that makes differential diagnosis explicit and optimizable; (2) three-dimensional pedagogical scoring integrated into GRPO optimization; and (3) a cognitive-role-driven multi-agent orchestration paradigm. Experiments on MIMIC-CXR and IU X-Ray validate the effectiveness of each component and the framework's overall superiority in both clinical accuracy and pedagogical readability.

%% ============================================================
%% ACKNOWLEDGMENTS
%% ============================================================
\subsection*{Acknowledgments}
--

%% ============================================================
%% BIBLIOGRAPHY
%% ============================================================
\begin{thebibliography}{21}

\bibitem{r2gen} Chen, Z., et al. Generating Radiology Reports via Memory-driven Transformer. \emph{EMNLP}, 2020.

\bibitem{radialog} Pellegrini, C., et al. RaDialog: A Large Vision-Language Model for Radiology Report Generation and Conversational Assistance. \emph{arXiv:2311.18681}, 2023.

\bibitem{tanno2023} Tanno, R., et al. Consensus, Dissensus and Synergy Between Clinicians and Specialist Foundation Models in Radiology Report Generation. \emph{arXiv:2311.18260}, 2023.

\bibitem{cxrclip} You, K., et al. CXR-CLIP: Toward Large Scale Chest X-ray Language-Image Pre-training. \emph{arXiv:2310.13292}, 2023.

\bibitem{imitate} Liu, C., et al. IMITATE: Clinical Prior Guided Hierarchical Vision-Language Pre-training. \emph{arXiv:2310.07355}, 2023.

\bibitem{chexagent} Chen, Z., et al. CheXagent: A Vision-Language Foundation Model to Enhance Efficiency of Chest X-ray Interpretation. \emph{arXiv:2401.12208}, 2024.

\bibitem{diseaseaware} Wu, P., et al. A Disease-Aware Dual-Stage Framework for Chest X-ray Report Generation. \emph{AAAI}, 2026.

\bibitem{s2dalign} Gao, J., et al. S2D-ALIGN: Shallow-to-Deep Auxiliary Learning for Anatomically-Grounded Radiology Report Generation. \emph{arXiv:2511.11066}, 2025.

\bibitem{cxrmate2} Nicolson, A., et al. CXRMate-2: Structured Multimodal Temporal Embeddings and Tractable Reinforcement Learning for Clinically Acceptable Chest X-ray Radiology Report Generation. \emph{arXiv:2604.18967}, 2026.

\bibitem{riha} Chen, Y., et al. RIHA: Report-Image Hierarchical Alignment for Radiology Report Generation. \emph{JBHI}, 2026.

\bibitem{marlrad} Baba, K., et al. MARL-Rad: Multi-Modal Multi-Agent Reinforcement Learning for Radiology Report Generation. \emph{arXiv:2603.16876}, 2026.

\bibitem{xrayclaw} Young, S., Xu, L. XrayClaw: Cooperative-Competitive Multi-Agent Alignment for Trustworthy Chest X-ray Diagnosis. \emph{arXiv:2604.02695}, 2026.

\bibitem{autogen} Wu, Q., et al. AutoGen: Enabling Next-Gen LLM Applications via Multi-Agent Conversation. \emph{arXiv:2308.08155}, 2023.

\bibitem{chexmix} Kumar, A., et al. CheXmix: Unified Generative Pretraining for Vision Language Models in Medical Imaging. \emph{CVPR Findings}, 2026.

\bibitem{maarta} Awasthi, A., et al. MAARTA: Multi-Agentic Adaptive Radiology Teaching Assistant. \emph{MICCAI}, 2025.

\bibitem{imactcxr} Le, T.A., et al. IMACT-CXR: An Interactive Multi-Agent Conversational Tutoring System for Chest X-Ray Interpretation. \emph{ISBI}, 2026.

\bibitem{cemrag} Salm\`{e}, M., et al. Concept-Enhanced Multimodal RAG: Towards Interpretable and Accurate Radiology Report Generation. \emph{arXiv:2602.15650}, 2026.

\bibitem{green} Ostmeier, S., et al. GREEN: A Clinical Efficacy Metric for Radiology Report Generation. \emph{NAACL}, 2024.

\bibitem{radgraph} Jain, S., et al. RadGraph: Extracting Clinical Entities and Relations from Radiology Reports. \emph{NeurIPS Datasets and Benchmarks}, 2021.

\bibitem{mimiccxr} Johnson, A.E.W., et al. MIMIC-CXR: A Large Publicly Available Database of Labeled Chest Radiographs. \emph{Scientific Data}, 2019.

\bibitem{chexpert} Irvin, J., et al. CheXpert: A Large Chest Radiograph Dataset with Uncertainty Labels and Expert Comparison. \emph{AAAI}, 2019.

\end{thebibliography}

\end{document}
